\section{Nearer Than Nearness}

\subsection{Creation has no material underneath it}

Human making begins with something: stone, sound, language, a field of forces, an inherited world. Creation in the strict sense does not. To create \latin{ex nihilo} is not to shape a mysterious substance called ``nothing.'' It is to cause the whole being of the creature without presupposing matter, subject, place, or time as an independent partner (ST I, q.~45, aa.~1--3).

This is why creation is not merely a very great change. Change presupposes a subject that becomes otherwise. In creation there is no prior creature undergoing the transition from nonbeing to being. There is only the creature as wholly caused, with everything by which it is something rather than nothing received from the universal cause. Even time belongs on the side of the gift.

St.~Athanasius begins his account of the Incarnation here. Against schemes in which the world emerges from pre-existent matter or from another creator, he confesses that the good God made the universe from nothing through his Word. The same Word who gives existence is the one who enters creation to rescue it from corruption (\emph{On the Incarnation} 3--5). Creation and redemption already lean toward one another: the source of being is not indifferent to what he has freely called forth.

``Freely'' is essential. The world is not a necessary spill from the divine substance, and God does not create to complete a lack. He necessarily loves his own infinite goodness; he does not need a finite reflection in order to be good. Creatures are willed without compulsion, not because the Creator is lonely, improved, or enlarged by them (ST I, q.~19, aa.~2--4; Vatican I, \emph{Dei Filius}, ch.~1).

\subsection{The gift is being given now}

Creation's radical character does not recede with the age of the universe. If God caused only an initial state and then withdrew, creatures would need to acquire independent possession of the act of being. But what has existence by participation never turns that reception into existence by essence. The dependence belongs to what the creature is, not merely to how long ago it began.

Aquinas therefore treats conservation as the continuation of the creative giving of being. It is not a repeated divine decision or a repair performed from moment to moment. God's one eternal act has a temporal effect: for as long as a creature exists, its act of existence remains caused. Aquinas compares this to light remaining in air while illumination continues, while warning that the divine giving itself implies neither motion nor succession in God (ST I, q.~104, a.~1).

The world is not a machine abandoned by its maker. Nor is it a stage on which God occasionally appears. The Creator is present to every creature as the cause of its most intimate actuality. Being is deeper in a thing than any color, location, feeling, memory, or relation; without it none of these can be actual. Because God causes that act, Aquinas says God is in all things ``most intimately''---not as a part of their essence or an accident within them, but as the agent giving them being (ST I, q.~8, a.~1).

Augustine found language for the paradox: God is \latin{interior intimo meo et superior summo meo}---more inward than what is innermost in me and higher than what is highest in me (\emph{Confessions} III.6.11). Divine transcendence and divine intimacy do not compete. Precisely because God is not one being at a distance from another, spatial separation does not mediate his causal presence. He is beyond every creature by nature and present at the root of its actuality as sustaining cause.

\subsection{Participation without absorption}

Aquinas uses participation to name this received likeness. Iron made hot by fire is genuinely hot, but it receives heat rather than possessing it as source. The analogy is limited: fire too is a creature whose being is received. Yet in the relevant order it shows how a perfection can truly belong to a participant without making that participant its underived source. So every creature genuinely exists and is genuinely good, but under a bounded mode it does not originate. What belongs to God by essence is reflected in creatures by participation (ST I, q.~44, a.~1).

The analogy protects two truths. Against pantheism, the creature is not a piece, mode, or dilution of God. Creation produces a reality really distinct from its cause. Against deism, that distinction is not independence. The creature is most itself when it receives; dependence upon the Creator is not an injury later imposed on its freedom but the condition under which there is a creature capable of freedom at all.

Look again at the hand, the wall, the small pool of light. Their familiarity has returned, but it has changed. Each is poised between nothing and the One who is, not precariously as though divine attention might wander, but gratuitously: nothing in the creature makes the gift owed. Its existence is a fact inside the world and, at the same time, a relation to the uncreated source no instrument can measure.

The metaphysical regress has therefore turned inward. The source is not only ``back there'' at the end of reasoning. The act by which the reader now exists is already the place where the dependence is enacted. Beneath thought, beneath breath, beneath the capacity to ask why, there is the silent, continuous gift: \emph{you do not have being from yourself. At this instant, you are being given to be.}
